\section{Related Work}
\label{sec:related_work}

\textbf{Static Parameter-Space Model Merging:} Model-merging methods aim to consolidate multiple specialized neural networks into a single cohesive model without undergoing expensive retraining. Early techniques like Model Soups \cite{wortsman2022model} and Git Re-Basin \cite{ainsworth2022git} average the weights of models that share a pre-trained initialization. To handle more diverse downstream tasks, Task Arithmetic \cite{ilharco2022editing} subtracts the pre-trained weights from fine-tuned weights to create task vectors, combining them via linear addition. TIES-Merging \cite{yadav2023ties} improves on this by resolving sign conflicts and pruning redundant parameters, while DARE \cite{yu2023dare} employs randomized delta pruning to preserve model capabilities. ZipIt! \cite{stoica2023zipit} merges models with disparate features by aligning redundant internal activations. While highly parameter-efficient and requiring zero runtime memory or compute overhead, these static methods suffer from ``heterogeneity collapse'' under realistic multi-task input streams. Because the parameters are merged statically, the model cannot dynamically adjust its routing on a per-sample basis, resulting in catastrophic cross-domain interference and degraded classification performance on interleaved streams.

\textbf{Dynamic PEFT Serving and Systems Pipelines:} Parameter-Efficient Fine-Tuning (PEFT) adapters, such as LoRA \cite{hu2021lora}, prefix-tuning \cite{li2021prefix}, and prompt tuning \cite{lester2021power}, isolate downstream task adjustments to a tiny subset of parameters, keeping the massive pre-trained base model frozen. To serve multiple concurrent LoRA adapters in production, high-throughput systems like S-LoRA \cite{sheng2023slora} and Punica \cite{chen2023punica} manage weight tensors using specialized CUDA paging and memory-management kernels on server-class GPU clusters. At the edge, where GPU clusters are unavailable and execution happens on resource-constrained CPUs or microcontrollers, such systems are impractical due to heavy scheduling and paging overheads. Instead, edge architectures utilize dynamic routing networks like PFSR \cite{pfsr2025} to route inputs to specialized adapters. To avoid cross-task interference under mixed-task streams, state-of-the-art pipelines employ Micro-Batch Homogenization (MBH) \cite{mbh2025} to partition heterogeneous batches into sequential task-homogeneous sub-batches. While effective at preserving accuracy, MBH forces the heavy base model backbone to run sequentially $G$ times (where $G$ is the number of active tasks in the batch), multiplying latency and DRAM-to-SRAM weight swapping overheads, and violating on-device real-time constraints \cite{chatterjee2024robustness}.

\textbf{Activation-Space Blending and SPS-ZCA:} To bypass the latency-heterogeneity trade-off, recent work has explored dynamic activation-space ensembling. SABLE (Sample-wise Activation Blending of Low-Rank Experts) \cite{huang2024lorahub} blends adapter activations layer-by-layer rather than merging weights, executing serving in a single forward pass. SPS-ZCA (Single-Pass Sample-Wise Routing with Zero-Shot Centroid Alignment) builds upon this by routing inputs at Layer 3 using head-independent representation task centroids pre-computed from a small calibration split. Since ensembling occurs on-the-fly via a parallel forward pass of the base model, SPS-ZCA achieves absolute immunity to heterogeneity collapse with a constant $O(1)$ backbone execution latency. However, existing activation-blending frameworks are restricted to high-precision floating-point execution (FP32/FP16), failing to leverage edge-native low-power integer accelerators.

\textbf{Quantization and Edge Serve Constraints:} Quantization is a foundational technique to make deep neural networks executable on edge microprocessors and NPUs by discretizing continuous parameters into low-bitwidth integers (INT8 or INT4) \cite{warden2019tinyml, dettmers2023qlora}. In modern multi-expert registries, quantizing independent experts introduces severe representation noise because different experts develop highly disparate dynamic ranges and activation outlier distributions. Applying naive quantization to the expert weights and activations destroys the delicate geometric alignment of the low-rank subspaces, resulting in catastrophic precision loss and accuracy degradation. Our proposed Q-SPS is the first to execute activation-space dynamic blending in pure integer precision. We introduce Quantization-Aware Scale Calibration (QASC) as a training-free structural defense that dynamically calibrates activation scaling factors, neutralizing low-bitwidth representation noise and enabling highly efficient edge-native execution.